% Triangle Mutiny — Stability Envelope (v1.2)
% Author: Borealis S. Hedling

\documentclass[11pt]{article}
\usepackage{amsmath, amssymb, amsthm}

\title{Triangle Mutiny: Stability Envelope (v1.2)}
\author{Borealis S. Hedling}
\date{October 2026}

\begin{document}
\maketitle

\section{Overview}

The stability envelope defines the region of the reasoning manifold where
reasoning remains deterministic, invariant-preserving, and collapse-free.
It integrates curvature continuity, spectral mismatch, modal adjacency,
and drift accumulation into a unified geometric constraint.

This module depends on:
\begin{itemize}
  \item \texttt{core-math.tex}
  \item \texttt{collapse-geometry.tex}
  \item \texttt{spectral-operators.tex}
  \item \texttt{modal-adjacency.tex}
\end{itemize}

\section{Holonomy Drift Tensor}

Define the holonomy drift tensor at point $p$:
\[
H(p)
  = H_{\mathrm{curv}}(p)
  + H_{\mathrm{spec}}(p)
  + H_{\mathrm{modal}}(p)
  + H_{\mathrm{drift}}(p),
\]
where:
\[
H_{\mathrm{curv}}(p) = |K(p^+) - K(p^-)|,
\]
\[
H_{\mathrm{spec}}(p)
  = \delta_{\mathrm{spec}}(p)
    + I_{\mathrm{spec-curv}}(p),
\]
\[
H_{\mathrm{modal}}(p)
  = d_{\mathrm{modal}}(\mu(p^-), \mu(p^+)),
\]
\[
H_{\mathrm{drift}}(p)
  = \|J(p)\|,
\]
and $J(p)$ is the Jacobi failure magnitude.

\section{Stability Envelope Definition}

The stability envelope is defined as:
\[
S = \{ p \in M : H(p) < \theta \text{ and } p \notin \Sigma \}.
\]

Points inside $S$ are stable; points outside are near-collapse.

\section{Stability Inequalities}

A point $p$ is stable when all inequalities hold:
\[
|K(p^+) - K(p^-)| < \theta_{\mathrm{curv}},
\]
\[
\delta_{\mathrm{spec}}(p)
  + I_{\mathrm{spec-curv}}(p)
  < \theta_{\mathrm{spec}},
\]
\[
d_{\mathrm{modal}}(\mu(p^-), \mu(p^+))
  < \theta_{\mathrm{modal}},
\]
\[
\|J(p)\| < \theta_{\mathrm{drift}}.
\]

\section{Stability Neighborhood}

Define the stability neighborhood:
\[
U(S)
  = \{ q \in M : \mathrm{dist}(q, S) < \delta \}.
\]

Points in $U(S)$ are monitored for drift accumulation.

\section{Stability Gradient}

Define the stability gradient:
\[
\nabla S(p)
  = \left(
      \frac{\partial H_{\mathrm{curv}}}{\partial p},
      \frac{\partial H_{\mathrm{spec}}}{\partial p},
      \frac{\partial H_{\mathrm{modal}}}{\partial p},
      \frac{\partial H_{\mathrm{drift}}}{\partial p}
    \right).
\]

The gradient determines how quickly a point approaches collapse.

\section{Stability Classification}

A point $p$ is classified as:
\begin{itemize}
  \item \textbf{stable} if $p \in S$,
  \item \textbf{near-collapse} if $p \in U(S)$,
  \item \textbf{collapsed} if $p \in \Sigma$.
\end{itemize}

\section{Deterministic Kernel Constraint}

The deterministic kernel enforces:
\[
p_{t+1} \in S.
\]

If $p_{t+1} \notin S$, the kernel triggers a governance correction.

\end{document}
